%% clawmarc_wadl.tex — reviewer-facing prong for WADL 2026 (ACM sigconf).
%% Sole author: A. Kornai; AI assistance in Acknowledgments (ACM policy).
%% Parallel file: clawmarc.tex (clawxiv-ready prong, joint authorship).
%% The two are synced BY HAND; never refactor into shared \input bodies.
%% v1.rc2 (2026-09-07): restructured after the referee-report draft; DCMI classes; pilot.
%% v1.rc3 (2026-09-11): rewritten after the adversarial review of 2026-09-11 —
%%   ACM reference strip restored; bundle/repository state stated honestly;
%%   model root cards; trust model and test-key labelling; manifest-binding
%%   language corrected; fixed-size claims marked as hypotheses; results table
%%   with query counts, proxy-MAP named as such; reproducibility declaration;
%%   comparison table; XLM-R and the four checkpoints cited exactly.
\documentclass[sigconf]{acmart}

\setcopyright{none}
\settopmatter{printacmref=true}
\acmDOI{}
\acmISBN{}
\renewcommand\footnotetextcopyrightpermission[1]{}
\acmConference[WADL 2026]{Web Archiving and Digital Libraries Workshop}{October 16, 2026}{Dallas, TX, USA}
\acmYear{2026}
\copyrightyear{2026}

\usepackage{booktabs}

\begin{document}

\title{The 4096-byte catalog card: fixed-size, signed, content-bound description for archives that machines read}

\author{Andr\'as Kornai}
\orcid{0000-0001-6078-6840}
\email{kornai@sztaki.hu}
\affiliation{%
  \institution{SZTAKI Institute for Computer Science and Control}
  \city{Budapest}
  \country{Hungary}
}

\begin{abstract}
We describe clawmarc, a catalog card for digital artefacts that is exactly
4096 bytes long, binds by SHA-256 to the bytes it describes, carries a short
human-readable description alongside a machine-readable search payload,
records retrieval hints for several storage backends, and carries an Ed25519
signature. The card is a descriptor, not a proof of authenticity or
availability, and it is meant to live separate from the artefact; its
artefact classes are the twelve terms of the DCMI Type Vocabulary, and subject
description is left to the schemes a cataloger already uses. We give the
design and its rationale, state the trust model under which a signature means
anything, place the card against MARC, Dublin Core, Linked Data, WARC,
Memento, and BagIt, and report what has been built: a specification with a C
reference layout, a Python producer and inspector, reference cards, cards
that pin the exact embedding models the format names, a differential test in
which two independent implementations produced byte-identical cards on every
pinned field, and retrieval experiments on a judged collection and a
topic-proxy collection. The experiments do not propose a retrieval method;
they ask whether a fixed card can carry representations from models with
materially different behaviour while keeping one catalog interface. The
contribution claimed is a substrate, and its costs and open questions are
stated with it.
\end{abstract}

\begin{CCSXML}
<ccs2012>
<concept>
<concept_id>10002951.10003317.10003318</concept_id>
<concept_desc>Information systems~Document representation</concept_desc>
<concept_significance>500</concept_significance>
</concept>
<concept>
<concept_id>10002951.10003317.10003338</concept_id>
<concept_desc>Information systems~Retrieval models and ranking</concept_desc>
<concept_significance>300</concept_significance>
</concept>
<concept>
<concept_id>10002951.10003227.10003392</concept_id>
<concept_desc>Information systems~Digital libraries and archives</concept_desc>
<concept_significance>500</concept_significance>
</concept>
</ccs2012>
\end{CCSXML}
\ccsdesc[500]{Information systems~Digital libraries and archives}
\ccsdesc[500]{Information systems~Document representation}
\ccsdesc[300]{Information systems~Retrieval models and ranking}

\keywords{catalog card, content addressing, metadata, fixed-size records, embeddings, DCMI Type Vocabulary, web archiving, interoperability}

\maketitle

\section{Introduction}

For about a century the unit of the library catalog was a card of
$7.5 \times 12.5$ centimetres. Its virtues were not typographic. A drawer of
cards could be counted, its space planned, a card pulled without disturbing its
neighbours, and a set of cards from one library merged into the drawers of
another; the Library of Congress sold printed cards precisely because every
recipient's drawers were the same size. The card did not need to be as long as
the book; it needed to be long enough to say what the book was, where it stood,
and who had described it. Krajewski's history of the format traces how much
intellectual infrastructure rested on that one fixed dimension
\cite{Krajewski:2011}.

MARC replaced the card with a variable-length record \cite{Avram:1975}, and
every later interchange format, from Dublin Core \cite{Kunze:2007} to OAI-PMH
\cite{Lagoze:2001} to the graph serializations of Linked Data
\cite{Kellogg:2020}, followed it in letting the record grow to fit the
description. This was the right choice for records that skilled catalogers
write and databases hold. Its suitability is less clear for the situation now
facing digital libraries and web archives: artefacts produced in volumes that
no cataloging staff can describe, a growing share of them produced by, or with,
language models; storage replicated across mirrors, institutional
repositories, and content-addressed peer-to-peer networks
\cite{Benet:2014,Alam:2016}; and description increasingly performed by machines
for machines, with the human reading a card only after a search has already
ranked it.

This paper describes a catalog card designed for that situation. A clawmarc
card is exactly 4096 bytes. It binds to the bytes of the artefact it describes
by SHA-256, so that a later copy found anywhere can be checked against it. It
carries a short human-readable description (title, abstract, keywords,
classification, and as much body prefix as fits) and, in the same card, a
machine-readable payload, in the default profile a 768-dimensional document
vector. It records retrieval hints for HTTP, IPFS, IPNS, and Swarm, and it
carries an Ed25519 signature \cite{Josefsson:2017} whose meaning depends on
how the signing key was published, a point Section~\ref{sec:trust} takes up.
It is a descriptor, not a certificate of authenticity, availability, or
quality; it is intended to remain useful even when the artefact is no longer
available.

The format grew out of ClawXiv, an archival workflow for research produced
jointly by humans and language models \cite{Kornai:2026a}, where the need to
catalog machine-produced bundles at machine speed was felt early.
Section~\ref{sec:card} states the design. Section~\ref{sec:rationale} gives
the reasons for its main choices, including the ones that run against current
practice, and states the trust model. Section~\ref{sec:related} places the
card among the formats already in broad currency. Section~\ref{sec:impl}
reports what has been built and measured, with a reproducibility declaration,
and Section~\ref{sec:governance} lays out what remains, most of it
experimental.

\section{The card}
\label{sec:card}

\subsection{Layer model}

Three layers are kept apart. The \emph{card} is the 4096-byte descriptor. The
\emph{artefact} is the thing described. The \emph{catalog} is a collection of
cards together with indices, shards, heads, admission policy, and governance.
The format specifies only the first. A card can outlive its artefact;
authenticity and availability are resolved at the artefact and catalog layers.
This separation is what allows the card to be storage-neutral: cards and
artefacts may travel over local filesystems, HTTP mirrors, IPFS, Swarm, or an
institutional repository, and the card neither knows nor cares which.

\subsection{Layout}

The card is divided into a 1216-byte header of fixed-position fields, a
2816-byte \emph{arena} whose interpretation depends on the artefact class, and a
64-byte footer holding two CRC-32 checks. Table~\ref{tab:layout} summarizes the
header. Every multi-byte field lies on its natural alignment boundary; the
reference C structure compiles to exactly 4096 bytes with no padding and is
nevertheless declared packed, with compile-time assertions on the total size
and on the offsets of the fields most likely to drift, so that the wire layout
cannot silently change across compilers or across edits.

\begin{table}[t]
\caption{Header fields, in wire order (offsets in hex; the arena occupies
0x4c0--0xfbf, the footer 0xfc0--0xfff).}
\label{tab:layout}
\footnotesize
\begin{tabular}{@{}rrp{0.66\columnwidth}@{}}
\toprule
offset & size & field \\
\midrule
000 & 8   & magic \texttt{CXCC}, layout major/minor \\
008 & 8   & arena split, arena class, size class, flags \\
010 & 32  & card issued, work created, work revised, sequence \\
030 & 128 & schema, object, manifest, previous-card SHA-256 \\
0b0 & 96  & issuer public key, card signature (Ed25519) \\
110 & 132 & issuer card reference; access mode, encryption profile, key \\
194 & 88  & responsible ORCID; author, classification, URL counts; author digests \\
1ec & 8   & license identifier \\
1f4 & 272 & Swarm reference, IPFS CID, IPNS name, HTTP hint, locator-set digest \\
304 & 76  & embedding profile, four text lengths, text flags, text and embedding digests \\
350 & 32  & target card (bounded indirection) \\
370 & 336 & reserved (zero) \\
\bottomrule
\end{tabular}
\end{table}

\subsection{Bindings}
\label{sec:bindings}

Four digests tie the card to the world. The \emph{object} digest is the SHA-256
of the artefact bytes: the card says ``I describe these bytes,'' and any later
copy can be checked. The \emph{schema} digest identifies the normative layout
the producer followed; in every card published with this paper it is the
SHA-256 of the reference C header, so a consumer knows which layout it is
reading. The \emph{manifest} digest is optional: when the producer supplies a
manifest (a build record, a source listing, a migration record, or, for the
model cards of Section~\ref{sec:profiles}, a two-field record naming the
model and its archive) the digest binds it, and when no manifest is supplied
the field is zero and nothing further is bound. The \emph{previous-card}
digest forms a supersession chain when a card is reissued. None of these is a
proof of authenticity: a valid signature proves only that the holder of a key
made a signed statement about an object hash, and what the key stands for is
a question the card cannot answer by itself.

\subsection{The arena}

The arena is split at a producer-chosen offset into a machine payload and a
human text. For text artefacts the payload, when present, is a document vector
under a named \emph{embedding profile}. Three profiles are assigned:
profile~1 is the 384-dimensional English-language \nolinkurl{bge-small-en-v1.5}
checkpoint \cite{BAAI:2023,Xiao:2023}, stored as little-endian binary16 in
768 bytes; profiles~2 and~3 are the 768-dimensional multilingual
\nolinkurl{bert-base-multilingual-cased} \cite{Google:2019,Devlin:2019} and
\nolinkurl{paraphrase-multilingual-mpnet-base-v2} \cite{Reimers:2021,Reimers:2020},
1536 bytes in the same encoding, leaving 1280 bytes for text; profile~2 is
the default. Against profile~1 the trade is deliberate: half the
human-readable prefix buys a vector that serves the multilinguality of any
larger archive. A fourth model, the untuned \texttt{xlm-roberta-base}
\cite{FacebookAI:2019,Conneau:2020b}, was evaluated alongside these and is
archived with them for comparison; no profile is assigned to it, for the
reason Section~\ref{sec:impl} reports. The human text holds title, abstract,
keywords, classification, and whatever prefix of the body fits; a digest over
the full used text and a truncation flag make the cut explicit. The
classification field is free text and is where a producer writes whatever
subject scheme it uses: a Library of Congress class, a Dewey number, an arXiv
category, a schema.org type, or a JSON-LD fragment. For images the payload may
be a small visual reduction and the text a caption. For opaque artefacts the
split is zero and the text is a category descriptor; a producer must not
fabricate a document vector for bytes that contain no running text.

\subsection{Classes}
\label{sec:classes}

The one-byte \emph{arena class} names what kind of thing is described. It is a
type, not a subject: the distinction an operating system draws between text,
image, executable, and archive, not the distinction a classification scheme
draws between endocrinology and epidemiology. Rather than mint such a
vocabulary, the format reuses one the profession already maintains.
Values 16--27 are the twelve classes of the DCMI Type Vocabulary
\cite{DCMI:2020}, assigned in the alphabetical order of that vocabulary
(Table~\ref{tab:classes}), so that the URI of a class is recoverable from its
number. Value 0 means no class asserted. Values 1 and 2 are structural and
belong to the format: an \emph{indirect} card describes an artefact that is
itself a concatenation of cards, readable at 4096-byte stride, and a
\emph{doubly indirect} card describes a concatenation of indirect cards.
Consumers descend at most two levels; this is aggregation with a bounded
depth, and it is how a catalog shards a million cards into a few hundred
artefacts without inventing a container format. Values 3--15 are reserved to
the format and 28--255 are unassigned.

\begin{table}[t]
\caption{Arena class values. DCMI terms carry the URI prefix
\texttt{http://purl.org/dc/dcmitype/}.}
\label{tab:classes}
\footnotesize
\begin{tabular}{@{}rl@{\hspace{1.5em}}rl@{}}
\toprule
0 & no class asserted        & 16 & Collection \\
1 & indirect (cards)         & 17 & Dataset \\
2 & doubly indirect          & 18 & Event \\
3--15 & reserved (structural) & 19 & Image \\
28--255 & unassigned         & 20 & InteractiveResource \\
      &                      & 21 & MovingImage \\
      &                      & 22 & PhysicalObject \\
      &                      & 23 & Service \\
      &                      & 24 & Software \\
      &                      & 25 & Sound \\
      &                      & 26 & StillImage \\
      &                      & 27 & Text \\
\bottomrule
\end{tabular}
\end{table}

\section{Design rationale}
\label{sec:rationale}

\subsection{Fixed size}
\label{sec:fixed}

The fixed size is a design choice, and the case for it is a case about
catalog operations rather than about parser performance. The operations are
these. A shard of cards can be addressed by ordinal, memory-mapped, and
scanned at constant stride. A consumer can reject a record whose length is
not 4096 bytes before parsing any field, and a record-level integrity check
localizes damage to one slot. Ranking by a query vector is a matrix product
over the vector regions of a shard as it lies on disk. An archive that knows
its holdings knows its catalog's size to the byte, a million artefacts being
four gigabytes of cards whatever the artefacts are, so that replication,
sharding, and the decision to mirror another institution's cards are
decisions about a known quantity.

These are affordances, and the honest statement of what has been shown is
that the fixed size \emph{permits} them without a separate index. It has not
been shown that they improve a digital-library workload: variable-length
formats have well-engineered equivalents, an SQLite table, a Parquet column,
a JSON-lines file with an offset index, or a plain NumPy matrix of vectors
beside a record file, and a comparison of raw-card scanning against those on
a realistic catalog is an experiment not yet run (Section~\ref{sec:governance}).
What the fixed size buys without measurement is simpler to state: one
artefact type for the whole catalog, so that a mirror is a copy of a file,
and a validation rule that a consumer in any language can apply before it
understands a single field.

The card is not the complete record. A variable-length description can remain
in the artefact layer or in a manifest bound by digest. The card is a bounded
descriptor intended for replication and catalog-level operations, not a
finding aid. The choice of 4096 bytes follows a common storage granularity,
although the correspondence is platform-dependent; it holds the digests,
keys, signature, and locators the header needs, a 768-dimensional vector in
half precision, and over a kilobyte of prose. The drawbacks are plain:
verbose abstracts are truncated, and short artefacts waste most of the arena,
as they waste most of a 4~KB filesystem block.

\subsection{A machine payload beside the human text}

Cataloging work increasingly aims at high precision, often by combining
orthogonal descriptors. Suppose the interest is endometriosis in Peru. A
paper about uterine diseases among the Quechua could be highly relevant, but
will not surface on a keyword match. The machine payload adds a document
vector that carries such loose associations, broader and narrower alike, and
does so for queries not anticipated when the card was written. It is stored
on the card so that a catalog can rank without fetching the artefact and
without a central index. This is an affordance for retrieval systems, not a
claim that the card or any particular embedding profile solves information
retrieval.

The card does not help with quality. Tens of thousands of papers are
genuinely about quantum gravity; they are not false positives, and they vary
in quality by orders of magnitude. What the card can do is make the
description honest about its own provenance: the signing key is named; the
embedding profile is named; the text that was embedded is digested. The
extraction and reduction policy under which the vector was computed is
recorded in the card's JSON sidecar, and it is bound into the card only when
the producer supplies a manifest for the manifest digest to cover; a reader
of the bare card can verify what bytes it describes and which profile it
names, but cannot infer from the card alone that the sidecar's policy is the
policy that was used. Two producers describing the same bytes will differ in
their text and therefore in their vectors, and consumers compare vectors by
cosine similarity rather than byte identity. Section~\ref{sec:impl} reports
how far that variation reaches in practice.

\subsection{Embedding profiles and their obsolescence}
\label{sec:profiles}

Embedding models and their runtimes change over time. The format therefore
does not freeze a single model. The profile field is sixteen bits and names
model, dimension, and encoding together; the arena split is per card, so
moving from a 384-dimensional to a 768-dimensional profile costs only text, a
1024-dimensional half-precision vector fits the same way, and a binary or
8-bit quantization fits with room to spare. When a catalog moves to a new
profile it reissues cards along the supersession chain, and a consumer
holding cards under two profiles ranks each population separately, as it
would two indices. What the format fixes is where the vector lives and how it
is described, not which model made it.

Which model made it is nonetheless a question the catalog must be able to
answer exactly, and a colloquial model name does not: the same weights
circulate in several serializations, from several hubs, no two files
byte-identical. The format's answer is to make the model an artefact like
any other. The weights and tokenizer of a profile's model, in one canonical
serialization, are archived, and a card is minted for the archive: a
Software-class card that binds the exact bytes by digest, names the model in
its manifest, carries HTTP, IPFS, and Swarm locators, and is signed. Four
such cards accompany this paper, one for each model evaluated in
Section~\ref{sec:impl}; their identifiers are given in the reproducibility
declaration there. A catalog that archives the model beside the cards it
validates no longer depends on any distribution hub's longevity, and a
consumer asking which model profile~2 names gets an answer with the same
shape, and subject to the same checks, as any other artefact in the drawer.

Two properties must be kept distinct here. \emph{Model identity} is
content-bound: the card pins the bytes, and a consumer can verify that it
holds the same weights and tokenizer files the profile names. \emph{Vector
reproducibility} is weaker. The vector a consumer recomputes depends further
on pooling, chunk boundaries and weighting, tokenizer runtime behaviour, and
the floating-point environment, all of which the present profiles leave
partly to the producer; the differential test below measures the residual
this leaves within one environment and does not bound it across
environments. The model card fixes the first property. The second is the
subject of the drift experiment proposed in Section~\ref{sec:governance}, and
until it is run a recomputed vector should be compared by cosine, not
expected to match.

\subsection{Content binding and format migration}

A digital library normalizes formats: a PDF is migrated to PDF/A, a TIFF is
rewrapped, headers change. Every such act changes the bytes, and a card bound
to the old bytes no longer matches the new file. This is the intended
behaviour of a fixity binding, and it is exactly the property that lets a
reader notice that an archive has changed under them \cite{Lerner:2017,
Aturban:2019}; but it means the workflow must be stated. A migrated artefact
is a new artefact. It gets a new card, whose previous-card digest points at the
old one and whose manifest digest binds the migration record, so that the
chain carries the history and the old card remains true of the bytes it
described. This is the same discipline BagIt applies to a bag's payload
manifest \cite{Kunze:2018}, lifted from the package to the catalog.

\subsection{Publication posture}

A card is a publication act. Even when the artefact is withheld, a card may
disclose enough to identify, summarize, retrieve, classify, or compare it; a
document vector reveals gist-level information about the text it represents
and is not anonymization. The producer therefore runs under an explicit
content policy, shipped as an editable JSON profile, that says whether a text
prefix, a vector, a thumbnail, OCR, or an archive member listing may be
stored. The policy is recorded in the sidecar, and a producer that must be
able to prove which policy it applied places the policy in a manifest and
binds the manifest. A profile suited to a public science library allows all
of these; a profile for a directory of mixed-rights material allows almost
none and records hashes only. The same institutional, legal, and ethical
review that applies to a finding aid applies to a card, and the format is
built so that a producer can obey it.

\subsection{Trust, keys, and admission}
\label{sec:trust}

A signature is only as meaningful as the way its public key reached the
verifier, and the format is explicit about three cases.

The published reference artefacts, the two reference cards and the four
model cards, are signed with \emph{test keys}: each private key is derived
deterministically from a seed label that the sidecar records in the clear.
This is intentional. It lets two implementations byte-compare a signed card,
which is what the differential test requires, and it makes every published
card reproducible by anyone. It also means that anyone can derive the private
key and mint a differently-worded card under the same public key. The
signature on these cards therefore attests reproducibility and nothing else,
the sidecars say so in a \texttt{key\_kind} field, and no reader should treat
the reference issuer as an identity. A catalog that publishes cards for
consumption rather than for testing signs with a key generated at random and
kept private, the producer accepts such a key from a file, and the public key
is distributed through a channel the consumer already trusts: an
institution's web page over TLS, a repository's key list, or a card for the
key itself, which the header's issuer-card reference exists to point at.

Admission is a catalog decision, not a card property. A card whose signature
verifies under an unknown key is syntactically valid; whether it is admitted
to a catalog is a matter of that catalog's policy, which in the simplest case
is a list of accepted issuer keys. Key rotation is expressed with the
machinery already present: the issuer publishes a new key, reissues its cards
along the supersession chain, and the chain records the change. Revocation
is likewise a catalog-layer act, a statement that cards under a given key are
no longer admitted after a given date, and the format supplies the fields
such a statement needs without prescribing its transport. Finally, the
embedding profile numbers are assigned in the specification by its editor and
by agreement among the catalogs that will read them; who should hold that
role in the long run is the registry question of Section~\ref{sec:governance}.

\section{Relation to existing practice}
\label{sec:related}

Dublin Core \cite{Kunze:2007} is the closest relative: a deliberately small
element set, meant for machine harvesting and cross-domain use, whose Type
Vocabulary the card adopts outright. The card differs in being fixed-size,
signed, and bound to the bytes described; it can be read as a signed,
content-bound carrier for a Dublin-Core-sized description plus a vector.
OAI-PMH \cite{Lagoze:2001} moves such records between repositories on request;
a shard of cards is the same idea with the transport removed, so that a mirror
is a copy of a file.

The metadata community's direction of travel for twenty years has been toward
graphs: Linked Data, RDF, JSON-LD \cite{Kellogg:2020}, and the shared
vocabulary of schema.org \cite{Guha:2016}. The card is not a replacement for
that work. A graph is extensible, mergeable, and expressive; the card is not.
The card's proposed role relative to a graph is an envelope: a fixed-size,
signed, content-bound unit that can carry a JSON-LD fragment or a schema.org
type in its classification text, can be replicated without a triple store,
and can be ranked without being parsed. Signposting \cite{VanDeSompel:2015}
makes a related bargain for the live web, exposing a few typed links so that a
machine can navigate a scholarly object without understanding it; a card is a
signpost that survives offline.

In web archiving, WARC \cite{ISO:2017} is the record of capture, and the card
is not a substitute for it; a card can describe a WARC file, or a member of
one, as it can describe anything with bytes. InterPlanetary Wayback
\cite{Alam:2016} put WARC payloads on IPFS and is the nearest precedent for
content-addressed replication of archival holdings; a card carries the IPFS
CID and Swarm reference of what it describes and can serve as a signed,
portable index entry for such a store. Memento \cite{VanDeSompel:2013} answers
``what did this URI look like at time~$t$'' from a live archive; a card
answers the narrower question ``which bytes did this key describe, and where
might they be,'' and does so offline. The fixity work of Aturban, Alam,
Nelson, and Weigle \cite{Aturban:2019}, and the demonstration by Lerner, Kohno,
and Roesner that the archived web can be altered from the present
\cite{Lerner:2017}, motivate the object binding directly: a signed digest
issued at cataloging time and held outside the archive is what lets a later
reader notice the change. BagIt \cite{Kunze:2018} does this for a package;
the card does it for a catalog entry. Table~\ref{tab:compare} sets the
formats side by side on the properties this paper is about.

\begin{table}[t]
\caption{The card beside the formats it is most often compared with. ``Offline''
means a consumer holding only the record and the artefact can verify the
binding without a network. Extensibility is the property the card gives up.}
\label{tab:compare}
\footnotesize
\setlength{\tabcolsep}{2.2pt}
\begin{tabular}{@{}lp{0.85cm}p{0.95cm}p{0.85cm}p{1.15cm}p{0.85cm}p{0.75cm}@{}}
\toprule
 & extens\-ible & content-bound & signed & transport assumed & offline verify & vector rank \\
\midrule
DC / OAI-PMH & yes & no & no & HTTP harvest & no & no \\
JSON-LD & yes & no & option & graph store & no & no \\
Signposting & yes & no & no & live HTTP & no & no \\
WARC & yes & per record & no & file & yes & no \\
Memento & n/a & no & no & live HTTP & no & no \\
BagIt & yes & manifest & option & package & yes & no \\
clawmarc & no & SHA-256 & yes & none & yes & yes \\
\bottomrule
\end{tabular}
\end{table}

\section{Implementation and experience}
\label{sec:impl}

\paragraph{What is public, and which version.} The bundle this paper
describes, clawmarc v1.rc3, contains the specification as an Internet-Draft
in Markdown, XML, text, and HTML; the C header that is the normative layout;
the Python producer, inspector, and general-file cataloger; three editable
content policies; batch scripts that mint indirect and doubly indirect
collection cards over a directory tree; the two reference cards and the four
model cards with their JSON sidecars; the model archives' checksum manifest;
and the experiment scripts and result files named below. The GitHub
repository at \url{https://github.com/kornai/clawmarc} mirrors the bundle
from commit \texttt{76d8c2d} (2026-09-12) onward; the 2.2~GB model archives
are served from their locators, not committed. Should the two ever differ, the bundle
deposited with the paper, identified by the digests in the reproducibility
declaration, is the normative one.

\paragraph{Reference layout.} The header declares the structure packed and
asserts at compile time that it is 4096 bytes, that the arena begins at 0x4c0,
that the footer begins at 0xfc0, and that every multi-byte field sits on its
natural boundary. A producer in any language can be checked against it by
compiling one file.

\paragraph{Producer and inspector.} The Python producer detects type by
content, extracts text with version-pinned tools (pdftotext, pandoc, detex),
normalizes to Unicode NFC, computes the document vector when policy allows,
packs the card, signs it, and emits a JSON sidecar recording the extraction
and reduction policy and the exact text that was embedded. The inspector
validates length, magic, split, reserved-zero fields, both CRCs, and the
signature, and recomputes the object binding when the artefact is at hand.
The producer signs with the deterministic test seed unless a key file is
supplied, and labels the sidecar accordingly (Section~\ref{sec:trust}).

\paragraph{Differential test.} Two implementations were written independently
by two different language-model systems working from the same specification,
without sight of each other's code, and each minted a card for the same
artefact under the same test seed. The cards agreed byte for byte on every
field the specification pins: object and schema digests, class, split,
embedding profile, ORCID, license, chain and indirection fields, and title.
They disagreed on every field the specification had left to producer judgment:
which of several valid mirror locators to record, whether to populate work
timestamps from the bundle, and what to put in the abstract and keywords,
where one producer wrote a gist and the other scraped the source. The
embeddings told the same story in numbers. Two runs of one implementation
agreed to cosine 0.9999999999. The same text segments run through the other
implementation's pipeline agreed to 0.989, which measures how much the
construction of the embedding input was unpinned. The two cards as minted
agreed to 0.932, which measures how much the text differed. The specification
was tightened on the first two points; the third is the residual that
consumers must expect, and the reference cards are published with that
expectation stated.

\paragraph{Implementation demonstration.} Twenty PDFs were drawn at random
(seed 20260907) from the 3,295 in the author's paper archive and minted on a
laptop with text extraction, a document vector under profile~2, and a body
prefix under the public-library policy. All twenty produced valid cards of
class Text. Extraction took 0.08~s per file on average and minting 5.9~s,
most of it interpreter start and model load that a batch producer pays once.
The embedded text averaged 74,000 characters and is the entire extracted
body: the vector is computed over the full text, however long, while the
card's human-readable prefix is truncated at its 1280 bytes, as the design
expects. Over the resulting 80~KB of cards, and without access to the source
files, card-only ranking and keyword lookup ran as designed: a keyword scan
of the text regions for \emph{grammar} returned the three linguistics papers
and nothing else, and nearest neighbours by cosine grouped the linguistics
papers and the neural-computation papers as a reader would. This is a
demonstration that the operations work on the cards alone; it is not
evidence about retrieval quality, which the two collection experiments below
address.

\paragraph{Model comparison.} Four models were run over the same twenty
documents with the pipeline otherwise fixed: full extracted text, chunks at
the model's native window, token-count weighted mean, final L2
normalization. Throughput on the laptop was flat across the three
768-dimensional multilingual candidates, 10,900--11,700 tokens per second,
and 32,700 for the 384-dimensional profile~1. Everything else differed. The
untuned \texttt{xlm-roberta-base} collapsed: all 190 pairwise cosines fell
between 0.996 and 1.000, the entire ranking signal squeezed into the third
decimal place, and no profile is assigned to it. Multilingual BERT
(profile~2) spread the same pairs over 0.839--0.991; the English-only
sentence-tuned profile~1 model over 0.592--0.959; and the multilingual
sentence-tuned profile~3 model over 0.156--0.903. The two sentence-tuned
models, which share no architecture, training data, or language coverage,
ranked the 190 pairs in the closest accord of any pair of models (Spearman
0.90), the two untuned encoders agreed that well only with each other (0.88),
and tuned--untuned agreement ran at 0.80 and below. On the 1,400 Cranfield
documents the same ordering recurs at lower levels (pair Spearman 0.74 for
the two tuned models, 0.43--0.63 elsewhere). A reading consistent with these
figures is that the tuned models approximate a common similarity which the
untuned ones see through noise; on a sample of twenty documents and one
benchmark this is a hypothesis, not a finding. Quantization, by contrast,
costs nothing measurable: rounding the vectors to binary16, the card's
storage encoding, left rank order identical for all four models (Spearman
1.000, every nearest neighbour unchanged).

\paragraph{Judged retrieval experiment.} To separate the catalog affordance
from the quality of any one embedding, the same four models and two lexical
baselines were evaluated on the Cranfield collection: 1,400 documents, 225
queries, and the collection's graded relevance judgments. The evaluation used
the full extracted document text and cosine ranking for the vector models;
BM25 and sublinear TF--IDF used the same document and query texts.
Table~\ref{tab:results} gives the figures. The untuned multilingual encoders
performed far below the rest. These results are evidence about this
benchmark and this implementation, not a general ranking of embedding
models. The public documentation of the BGE and MPNet checkpoints does not
disclose Cranfield among their training resources \cite{BAAI:2023,Reimers:2021};
this is an absence of disclosed overlap, not proof that none exists. The
result supports the narrower conclusion that a card can carry
representations from models with materially different retrieval behaviour,
while leaving model selection to the retrieval application.

\begin{table}[t]
\caption{Retrieval results. Cranfield: 225 judged queries over 1,400
documents, graded relevance. NewReuters: 100 query documents over 5,000
sampled Reuters documents, with a shared topic label as the only relevance
proxy; the proxy metrics are not comparable with the judged ones and are
named accordingly. Single runs; no confidence intervals are available, and
differences of a few hundredths should not be read.}
\label{tab:results}
\footnotesize
\setlength{\tabcolsep}{2.4pt}
\begin{tabular}{@{}lrrrrrr@{}}
\toprule
 & \multicolumn{4}{c}{Cranfield (judged)} & \multicolumn{2}{c}{NewReuters (proxy)} \\
\cmidrule(lr){2-5}\cmidrule(lr){6-7}
model & MAP & P@10 & nDCG@10 & R@100 & proxy-MAP & proxy-P@10 \\
\midrule
BGE-small (p.\,1) & 0.329 & 0.254 & 0.359 & 0.758 & 0.500 & 0.730 \\
MPNet (p.\,3)     & 0.248 & 0.196 & 0.281 & 0.666 & 0.520 & 0.756 \\
mBERT (p.\,2)     & 0.073 & 0.054 & 0.078 & 0.349 & 0.456 & 0.579 \\
XLM-R (none)      & 0.019 & 0.015 & 0.020 & 0.126 & 0.369 & 0.441 \\
BM25                  & 0.274 & 0.226 & 0.315 & 0.697 & 0.442 & 0.736 \\
TF--IDF               & 0.259 & 0.205 & 0.295 & 0.699 & 0.440 & 0.711 \\
\bottomrule
\end{tabular}
\end{table}

\paragraph{NewReuters proxy experiment.} A second collection was built by
sampling 5,000 English Reuters documents and 100 query documents (seed
20260910), using shared Reuters topic labels as the only proxy for relevance.
The proxy-MAP and proxy-P@10 columns of Table~\ref{tab:results} report it.
These figures are not comparable with Cranfield MAP: the NewReuters queries
have no human relevance judgments, and the topic proxy can both miss relevant
documents and count merely topical documents as relevant. Within that
limitation, mBERT's proxy-MAP is competitive with the lexical baselines, but
its proxy-P@10 is well below theirs. The experiment therefore suggests, but
does not establish, a first-pass use for a broadly available multilingual
profile: candidate-set recall at a specified depth must be measured with real
judgments before deployment, and the Cranfield figures show that such a
profile should not be presented as a fine-grained final ranker. A catalog
that uses one profile to reduce the candidate set and delegates final ranking
to a task-specific index is a testable engineering hypothesis, not a result
of this experiment.

\begingroup\sloppy
\paragraph{Reproducibility declaration.} Layout~1.3; schema digest (SHA-256
of \nolinkurl{clawmarc_catalog_card.h}) \texttt{1dfaee46\ldots}, bound by
all six published cards. Reference cards \texttt{d71d5da6\ldots} (bundle
with a profile-1 vector) and \texttt{bf8a9b06\ldots} (general file, no
vector). Model cards (full names in Section~\ref{sec:profiles} and in the
bundle's manifest): BGE-small \texttt{1e6f3f41\ldots}; mBERT
\texttt{52e67acd\ldots}; MPNet \texttt{2071b96f\ldots}; XLM-R
\texttt{a452acf6\ldots}; each binds its archive by SHA-256, and the
checkpoint's \nolinkurl{model.safetensors} digest is recorded in the bundle's
model-identity note and in the bibliography entries. All six signatures are test-seed signatures. Runtime for every
measurement: Python 3.12.9, torch 2.14.0, transformers 4.49.0,
sentence-transformers 5.6.0, numpy 2.2.3, CPU inference in float32 with
binary16 storage, on an Apple M2 laptop (26~GB) under macOS~26.6. Chunking:
480 body tokens for profiles~1 and~2 and for XLM-R, 126 for profile~3 (its
native window); stride 0; token-count weighted mean; final L2
normalization. Seeds: 20260907 (demonstration and model comparison),
20260910 (both collection experiments). Result files, with SHA-256
prefixes: \nolinkurl{cranfield_experiment_results.json} \texttt{fd492638};
\nolinkurl{newreuters5000_results.json} \texttt{870b21da};
\nolinkurl{pilot_run_20260907_mbert_fulltext.json} \texttt{4c78183b};
\nolinkurl{model_compare_20260908.json} \texttt{a2b53220}; scripts
\nolinkurl{run_ir_experiments.py} \texttt{7feeeff9} and
\nolinkurl{run_cranfield_experiment.py} \texttt{cfca670d}. All are in the
bundle, whose manifest lists the full digests. The model archives have two
independent HTTPS sources: \url{https://lebadus.ai/Resources/Models/}, the
cards' HTTP hint, and the assets of the GitHub release \texttt{v1.rc3}, the
cards' second mirror; they are also on IPFS and Swarm. All four archives were
verified byte-exact against \nolinkurl{SHA256SUMS} on every route on
2026-09-12. The checksum manifest and the cards travel with the bundle and
the release, and are being mirrored to the lebadus endpoint.
\endgroup

\section{Future work}
\label{sec:governance}

The artefact classes now reuse a vocabulary the Dublin Core Metadata
Initiative maintains, and the structural values are fixed by the
specification, so the class byte needs no registrar. Two enumerations still
do: the embedding profiles, which will grow as models change, and the
encryption profiles. The specification leaves most of each space unassigned,
names no authority beyond its editor, and asks producers not to consume
unassigned values in public cards except by agreement among the catalogs
that will read them. Three resolutions are open: a registry lodged with an
existing body, profiles named by URI in a manifest and resolved rather than
enumerated, or continued enumeration by agreement while producers remain
few. The format has been written so that any of the three can be adopted
later without a layout change, and the choice will be made by the number of
producers rather than by argument.

\paragraph{Fixed size, measured.} The affordances of
Section~\ref{sec:fixed} should be costed against their variable-length
equivalents: raw-card stride scanning and matrix ranking against an SQLite
table with a vector column and against a NumPy matrix beside a JSON-lines
file, on a catalog of a million cards, measuring validation, random access,
ranking, and replication. Until that is done the fixed size rests on the
argument from uniformity, and the paper says so.

\paragraph{Model identity and drift.} The model cards make the identification
of a profile's model exact, but one quantity in them is recorded rather than
pinned: the runtime. A recomputed vector depends on the pooling code and the
floating-point environment as well as on the weights, and while the
differential test bounds reproducibility within one environment at ten
decimal places, drift across environments is unmeasured. Measuring it is the
next experiment, run the way the differential test was run: the same model
card, several environments, a cosine agreement matrix. If cross-environment
drift lands near the within-environment bound, digest-pinning the weights
suffices and the runtime stays advisory; if not, the profile must also name
the runtime.

\paragraph{Choice of model.} Which model deserves the next profile is an
empirical question involving retrieval quality, language coverage, runtime
cost, and the durability of the weight distribution. The two collection
experiments illustrate why the card should not prescribe that choice. BGE
and MPNet were the strongest vector profiles in the judged and proxy
evaluations respectively; mBERT was weak on the fine-grained Cranfield
judgments but competitive as a coarse pre-selector on the proxy; and XLM-R
was weak in both. The results do not establish a universal model ranking.
They establish that the profile field can preserve model identity while
catalogs choose different profiles for candidate generation and final
ranking.

\paragraph{Reduction.} The reduction from an unbounded text to one vector is
recorded in the sidecar but left to producer choice: chunk size, chunk
weighting, and pooling are all defensible in several ways, and the 0.989
agreement figure in the differential test measures exactly the cost of
leaving them unpinned. Candidate reductions are compared by the same offline
IR harness, and the specification will pin one when the comparison says
something.

\section{Conclusion}

The format described here assumes that the descriptor of a digital artefact
should have the same shape everywhere, and derives the rest: content binding
so the descriptor can be checked, a signature so it can be attributed once
its key is trusted, locators so it can be followed, a machine payload so it
can be ranked without being fetched, and a fixed size so that all of this can
be replicated and scanned as a dense array. Its purpose is to provide a
uniform substrate on which retrieval and cataloging methods can continue to
improve, while keeping older cards valid as models and storage systems
change. Cards minted for large public collections also provide durable
inventories against which later audits can detect bitrot, disappearance, or
loss of availability; the card itself does not restore a missing artefact.
The specification, reference implementation, reference cards, and model cards
are public in the bundle deposited with this paper, signed with keys that
are explicitly for testing. The measurements reported are preliminary and
single-run; the claims that rest on them are stated as hypotheses where they
are hypotheses, and the questions that require governance are separated from
the ones that require experiments.

\begin{acks}
We thank Eliot Lear (Cisco) for penetrating criticism of an earlier draft.
The specification and the reference implementation were developed with
substantive assistance from GPT-5 Codex (OpenAI) and Claude Opus (Anthropic),
including independent implementations of the differential test; Gemini 3.1
Pro (Google) ran an initial pilot that exposed a missing PDF branch in the
cataloger. GPT-5 Codex performed substantive experimental design,
implementation, and analysis for the collection experiments and wrote two
adversarial reviews, the second of which (2026-09-11) shaped this revision.
Claude Fable 5.1 (Anthropic) drafted an earlier version from the
specification and project record, minted the model cards, and revised the
paper in response to the reviews. The paper is submitted under the sole
authorship of Andr\'as Kornai, who is responsible for its content. This
disclosure follows ACM policy on authorship and generative AI.
\end{acks}

\bibliographystyle{ACM-Reference-Format}
\bibliography{clawmarc}

\end{document}
